\chapter{Технологический раздел}

\section{Выбор средств программной реализации}

В ходе реализации данного проекта был использован язык программирования С++ для разработки протокола. 
Такой выбор был сделан по следующим причинам:

\begin{itemize}
    \item C++ — компилируемый язык со статической типизацией, позволяющий получать высокопроизводительный машинный код и при необходимости работать с памятью и сетевыми буферами на низком уровне. 
    Это особенно важно при реализации бинарного протокола, где требуется явный контроль над форматом сообщений, выравниванием и преобразованием типов.
    \item Для C++ существует большое количество готовых библиотек для работы с сетью (Boost.Asio, standalone Asio, Qt Network и др.), а также для сериализации данных. 
    Это позволяет сосредоточиться на логике протокола и уменьшить объём «ручного» кода для работы с сокетами и обработкой соединений.
    \item Современные стандарты C++ (C++17/C++20) предоставляют инструменты для асинхронного и конкурентного программирования: потоки, futures, std::async, корутины. 
    Это позволяет реализовать неблокирующую обработку клиентских подключений и сообщений, что важно для протокола совместного редактирования, обслуживающего сразу несколько пользователей.
\end{itemize}


Для реализации сетевой части протокола совместного редактирования заметок в проекте использованы современные средства асинхронного программирования на языке C++. 
Такой выбор обусловлен необходимостью обработки сразу нескольких клиентских подключений без блокировки основного потока и без потери отзывчивости системы. 
В частности, применяются следующие механизмы:

\begin{itemize}
    \item Асинхронная модель обработки событий;
    \item Корутины и механизм co\_await (при использовании C++20);
    \item Асинхронный реактор для обработки событий;
\end{itemize}

Для организации сетевого взаимодействия и реализации бинарного протокола совместного редактирования текстовых заметок в работе использовались низкоуровневые средства стандартного POSIX-интерфейса для работы с сокетами. В частности, применялись системные заголовочные файлы <sys/socket.h>, <netinet/in.h> и <arpa/inet.h>, обеспечивающие доступ к механизмам TCP/IP-взаимодействия. Использование данного подхода позволило реализовать полный контроль над процессом установления соединений, приёма и передачи бинарных данных, а также эффективно управлять сетевыми ресурсами на стороне сервера и клиента.

Для реализации криптографических механизмов защиты передаваемых данных в работе была использована библиотека OpenSSL, в частности её высокоуровневый интерфейс EVP, предоставляемый заголовочными файлами <openssl/evp.h> и <openssl/rand.h>. Данный интерфейс обеспечивает унифицированный доступ к современным алгоритмам симметричного шифрования, генерации криптографически стойких случайных значений и управлению контекстами шифрования. Применение OpenSSL позволило реализовать надёжное шифрование передаваемых сообщений протокола и обеспечить конфиденциальность данных при передаче по открытым сетевым каналам.

В ходе реализации проекта для отладки работы протокола был использован Wireshark [10]. Данный инструмент предназначен для детального анализа сетевых потоков данных посредством их визуализации.

\section{Структура проекта}

Проект представляет собой клиент-серверную систему для совместного редактирования текстовых заметок, взаимодействующую через разработанный бинарный протокол прикладного уровня. Передача данных осуществляется по TCP-соединению. Протокол поддерживает аутентификацию и регистрацию пользователей, получение списка заметок, создание и открытие заметок, передачу операций редактирования, синхронизацию версий документа, совместный доступ и разрешение конфликтов (включая механизм согласования объединения изменений).

Проект состоит из следующих основных модулей:

\begin{itemize}
    \item NoteServer - серверная часть системы, принимающая TCP-соединения клиентов, выполняющая обработку запросов, взаимодействующая с хранилищем и рассылающая обновления и уведомления участникам совместного редактирования;
    \item NoteClient - клиентское приложение, устанавливающее соединение с сервером, формирующее запросы протокола и обрабатывающее запросы/ответы;
    \item NoteProtocol - статическая библиотека, определяющая формат сообщений протокола и реализующая кодирование/декодирование запросов и ответов.
\end{itemize}

\subsection*{NoteServer}

Модуль NoteServer реализует серверную часть системы совместного редактирования. Сервер ожидает входящие TCP-соединения, обрабатывает запросы клиентов в соответствии с кодами операций протокола, взаимодействует с хранилищем (репозиторием) и обеспечивает рассылку уведомлений и обновлений всем участникам совместного редактирования.

Основные компоненты:

\begin{itemize}
    \item Сетевой слой: запуск TCP-сервера (run), принятие подключений, обработка клиентских соединений в отдельных потоках (handleClient, processClientMessages).
    \item Декодирование и диспетчеризация сообщений: определение типа операции и выбор соответствующего обработчика (processClientMessage).
    \item Механизм рассылки обновлений: отправка уведомлений всем подключенным клиентам (broadcastToAllClients) и адресная отправка ответов.
    \item Работа с хранилищем: взаимодействие с модулем Repository для проверки прав доступа, чтения/записи текста заметок, ведения версии и получения списка пользователей, имеющих доступ к заметке.
\end{itemize}


\subsection*{NoteClient}

Модуль NoteClient реализует клиентскую часть системы. Клиент устанавливает TCP-соединение с сервером, формирует запросы протокола, отправляет их и асинхронно обрабатывает ответы и уведомления от сервера. Клиент также инициирует операции редактирования заметок и участвует в процедурах синхронизации и разрешения конфликтов.

Основные компоненты:

\begin{itemize}
    \item Методы управления пользователем: sendAuthRequest, sendRegistrationRequest — авторизация и регистрация.
    \item Методы управления заметками: sendGetNotesRequest, sendCreateNoteRequest, sendOpenNoteRequest — получение списка, создание и открытие заметок.
    \item Методы редактирования и синхронизации: sendUpdateTextRequest, sendSyncNoteRequest — передача изменений и запрос синхронизации версии.
    \item Методы совместного доступа и конфликтов: sendShareNoteRequest, sendApproveMergeRequest, sendOwnerApproveMergeResponse — предоставление доступа, запрос на объединение изменений и подтверждение/отклонение merge владельцем.
    \item Обработка входящих сообщений (сервер → клиент): слот onReadyRead и обработчики handle* (например handleGetNotesResponse, handleUpdateTextMergedRequest и др.), выполняющие разбор пакетов и генерацию событий UI (сигналы).
\end{itemize}


\subsection*{NoteProtocol}

Модуль NoteProtocol реализует бинарный протокол совместного редактирования заметок. Он включает описание поддерживаемых операций и структуры сообщений, а также функции сериализации и десериализации данных для передачи между клиентом и сервером.

Основные компоненты:

\begin{itemize}
    \item Requests — набор структур запросов протокола (например: AuthRequest, RegistrationRequest, GetNotesRequest, CreateNoteRequest, SyncNoteRequest, OpenNoteRequest, UpdateTextRequest, ShareNoteRequest, ShareNoteNotifyRequest, ApproveMergeRequest, OwnerApproveMergeRequest, ServerApproveMergeRequest, UpdateTextMergedRequest);
    \item Responses — набор структур ответов и серверных уведомлений (например: AuthResponse, RegistrationResponse, GetNotesResponse, CreateNoteResponse, OpenNoteResponse, UpdateTextResponse, SyncNoteResponse, OwnerApproveMergeResponse);
    \item Protocol — модуль, реализующий методы кодирования и декодирования сообщений (например: encode*, decode*), а также определяющий коды операций (Operation) для идентификации типа пакета.
\end{itemize}



\section{Пример работы протокола}

На рисунке 3.1 показан пример работы протокола в виде последовательности пакетов, передаваемых по сети, с указанием времени их отправки, используемого протокола, длины пакета и дополнительной информации. На рисунке можно увидеть, что клиент отправляет данные серверу (длина пакета равна 419), после чего сервер отправляет клиенту пакет с кодом подтверждения получения данных (длина пакета равна 3).

Также на рисунке выделен бинарный пакет с данными редактирования, пояснение к которому приведено в таблице 3.1.

Данные примеры демонстрируют типичный сценарий взаимодействия клиента и сервера в рамках протокола совместного редактирования заметок.

\begin{figure}[H]
	\centering
	{\includegraphics[width = \textwidth]{../img/work_example.jpg}}
	\caption{Пример работы протокола в Wireshark}
	\label{fig:page-03}
\end{figure}


\begin{table}[H]
    \caption{Структура пакета совместного редактирования}
    \label{dump}
    \begin{center}
        \begin{tabular}{| p{5cm} | p{3.5cm} | p{4.5cm} |} 
            \hline
            \textbf{Поле} & \textbf{Длина (байты)} & \textbf{Значения в дампе} \\
            \hline
            OP\_TYPE & 1 & 01 \\
            \hline
            PAYLOAD\_LEN & 2 & 00 34 (52) \\
            \hline
            PAYLOAD (AES зашифрованные данные) & 52 & Блок данных \\
            \hline
        \end{tabular}
    \end{center}
\end{table}


\section{Тестирование системы}

В рамках тестирования системы совместного редактирования заметок была развернута тестовая среда, состоящая из виртуальных машин на базе операционной системы Ubuntu Linux. Была создана отдельная виртуальная машина для серверной части приложения, а также несколько виртуальных машин для клиентских приложений, которые были объединены в единую сеть для обеспечения взаимодействия.

На серверной виртуальной машине была развернута и запущена серверная часть приложения, отвечающая за обработку запросов клиентов, синхронизацию изменений документов и управление сессиями совместного редактирования. Клиентские виртуальные машины были настроены для подключения к серверу по указанному сетевому адресу и порту.

Тестирование основных функциональных возможностей системы проводилось через последовательную работу на клиентских машинах. Проверялись следующие сценарии:

\begin{enumerate}
    \item Одновременное подключение нескольких клиентов к серверу и создание общей сессии редактирования;
    \item Редактирование текстового документа с разных клиентских устройств в реальном времени;
    \item Синхронизация изменений между всеми подключенными клиентами;
    \item Обработка конфликтных ситуаций при одновременном внесении изменений в одну и ту же область документа;
    \item Корректное отображение состояния документа после переподключения клиента.
\end{enumerate}